\section{Odbiór, utrzymanie i podział prac}
\begingroup\small
\begin{tabularx}{\linewidth}{@{}P{24mm}Y Y@{}}
\toprule\textbf{Próba} & \textbf{Kryterium} & \textbf{Dowód odbioru}\\\midrule
CCTV dzień/noc & Kadr odpowiada zadaniu; czytelny ruch, brak istotnych zasłonięć i olśnień. & Próbki nagrań ze scenariuszy i zatwierdzona lista kadrów.\\
Zapis i odtworzenie & 24 kanały, wymagany czas retencji, kontrolowany eksport i bufor kart. & Raport obciążenia, próba odtworzenia i integralności eksportu.\\
KD i SSWiN & Prawidłowa akceptacja/odmowa, alarm otwarcia, sabotaż i praca bez LAN; W-01. & Macierz uprawnień i protokół z czasem obsługi.\\
SSP i ewakuacja & Próba wszystkich elementów i sterowań, potwierdzenie scenariusza, wyjście bez LAN. & Protokół branżowy; pomiary sygnalizacji; zgody i dokumenty urządzeń.\\
LAN i okablowanie & Certyfikacja miedzi, pomiary włókien, odtworzenie po przerwaniu R1--R4; W-02. & Raporty pomiarów i czasy rekonwergencji dla każdego odcinka.\\
Zasilanie & 30 min UPS przy obciążeniu, zanik sieci, start agregatu i test baterii DC. & Rejestr mocy, napięć i czasu; próba SZR.\\
Awaria integracji & Brak VMS/AI/LAN nie wstrzymuje lokalnych funkcji SSP; W-03. & Scenariusz awaryjny i podpisy branż odpowiedzialnych.\\
AI i dane & Raport TP/FP/FN, czas p95, obciążenie operatora, maski i role. & Niezależny zestaw testowy i raport bez deklaracji niezmierzonych wyników.\\\bottomrule
\end{tabularx}
\endgroup

\subsection*{Utrzymanie}
Codziennie: przegląd alarmów, zapisu i zasilania przez dyżurnego. Co miesiąc: proponowana próba wybranego scenariusza i kontrola kopii. Kwartalnie: proponowana próba odtworzenia i analiza jakości alarmów. Terminy branżowych przeglądów wynikają z właściwych przepisów i instrukcji producentów; dla urządzeń ppoż. rozporządzenie przewiduje przeglądy nie rzadziej niż raz w roku [2].

\subsection*{Propozycja pracy zespołowej -- według pakietów}
\begin{enumerate}
\item \textbf{Osoba A:} obiekt, ryzyka, geometria i projekt trasowy.
\item \textbf{Osoba B:} CCTV, KD i SSWiN; rejestr urządzeń.
\item \textbf{Osoba C:} SSP, scenariusze reakcji i zasilanie.
\item \textbf{Osoba D:} sieć, AI, zestawienie i redakcja PDF. Przy innym składzie pakiety można łączyć.
\end{enumerate}
Na prezentacji warto przejść przez jeden alarm od urządzenia do reakcji operatora, następnie pokazać tę samą sytuację przy braku LAN i energii. Taka narracja łączy schematy, obliczenia i procedury.
